% ============================================
% MEM-21
% Gerado de molambudos.md — build_miolo.py
% ============================================
\fragtitle{mem}{M}{MEM}{21}{O Jornal}
\markboth{MEM-21 --- O Jornal}{}

\noindent \textit{1948 --- a notícia que mentia}

\noindent \textbf{1948 --- Uma notícia do lado de fora}

\vspace{0.5em}\noindent\rule{\textwidth}{0.5pt}\vspace{0.5em}

\noindent Um enfermeiro deixou um jornal na cela. Não sei por quê. Talvez por esquecimento. Talvez por pena. O jornal era de Belo Horizonte, de 1947. Tinha cheiro de mofo e de fumaça de cigarro alheio. A textura das páginas era quebradiça, amarelada como folha seca --- eu segurei com cuidado, como se tudo que eu ainda acreditava sobre o mundo lá fora também pudesse desmanchar entre os dedos. Uma notícia na página 5 me fez parar:

\noindent \textit{\textquotedblleft{}COLÔNIA DE BARBACENA COMPLETA 40 ANOS COM 5 MIL INTERNADOS\textquotedblright{}}

\noindent Eu li a matéria. Dizia que o Colônia era \textquotedblleft{}um dos maiores hospitais psiquiátricos da América Latina\textquotedblright{}. Dizia que o governo investia em \textquotedblleft{}terapias modernas\textquotedblright{}. Dizia que os pacientes \textquotedblleft{}recebiam tratamento humanizado\textquotedblright{}.

\noindent Mentira.

\noindent Tudo mentira.

\noindent Eu estava ali. Eu via os eletrochoques sem anestesia. Eu via os pacientes morrerem de pneumonia sem remédio. Eu via os corpos empilhados na carroceria do caminhão todas as manhãs.

\noindent O jornal mentia.

\noindent E eu pensei: \textit{\textquotedblleft{}Se o jornal mente sobre isso, sobre o que mais ele mente?\textquotedblright{}}

\noindent A criatura respondeu:

\noindent \textit{\textquotedblleft{}Ele mente sobre tudo. Os jornais sempre mentem. Eles dizem que o mundo é organizado. Não é. O mundo é uma vala. E a vala está cheia de corpos que o governo não quer que você veja.\textquotedblright{}}

\noindent \textit{\textquotedblleft{}Por que você está me dizendo isso?\textquotedblright{}}

\noindent \textit{\textquotedblleft{}Porque você é a prova. Você está vivo. Você viu. Você pode contar.\textquotedblright{}}

\noindent \textit{\textquotedblleft{}Quem vai acreditar em mim?\textquotedblright{}}

\noindent \textit{\textquotedblleft{}Quem ler o diário.\textquotedblright{}}

\noindent Ela estava certa.

\noindent O jornal mentiu.
O diário não.

\noindent O diário é a verdade que o jornal nunca publicou.

\noindent Guardei aquela página do jornal dentro do colchão, junto com os primeiros cadernos do diário. Não para provar que ele mentiu. Para lembrar, sempre que eu duvidasse de mim mesmo, que eu não inventei nada do que vi.

\vspace{0.5em}
\noindent\textit{\textbf{↪ Rotas:} \rota{MEM-22} \rotasep \rota{DOC-19}}